\section{HPS Apparatus, Data-Taking Campaigns, and Analysis Scope}
\label{sec:datasets}

\subsection{Detector configurations and apparatus relevant to the prompt search}

HPS is a forward spectrometer optimized for narrow $e^+e^-$ resonances and displaced
vertices in the high-rate environment immediately downstream of a thin tungsten target.
Its prompt $A^\prime \rightarrow e^+e^-$ sensitivity is set primarily by three
subsystems: the Silicon Vertex Tracker (SVT), which determines the mass and vertex
resolution; the electromagnetic calorimeter (ECal), which provides the trigger and
track--cluster matching; and the beamline/target geometry, which defines the forward
acceptance and the occupancy environment \cite{HPSExperiment2022,HPS2016PromptLong}.

The 2015 and 2016 engineering runs used the baseline six-station SVT. In that
configuration, six axial/stereo tracking stations were distributed between
$z=\SIrange{10}{90}{cm}$ downstream of the target, with the first active silicon only
about \SI{1.5}{mm} from the beam plane \cite{HPS2016BumpInternal,HPSExperiment2022}.
The ECal geometry used in the engineering runs was already close to the final HPS
layout: two PbWO$_4$ crystal halves with the highest-rate crystals nearest the beam
removed, leaving the beam to pass between the two halves
\cite{HPS2016BumpInternal,HPSExperiment2022}.

Before the first physics run in 2019, HPS upgraded the front tracker. The 2019 and
2021 running periods used a seventh double layer at $z=\SI{5}{cm}$, thin silicon,
slim-edge processing, and short split striplets to preserve low-mass acceptance while
reducing material and occupancy \cite{HPSExperiment2022,HPSSlimEdgeSensors2020}. The
same slim-edge module design was used to replace the previous first layer, and the
front layers were moved closer to the beam. Those hardware changes are the main reason
the 2019/2021 configuration improves the low-mass prompt search: the first precision
measurements occur closer to the beam, the front material budget is reduced, and the
small-opening-angle acceptance is correspondingly better. The ECal crystal layout
itself did not change between the engineering and physics runs, but a positron
hodoscope was installed in front of the ECal before the 2019 run to support the
single-arm positron trigger that recovers events otherwise lost in the ECal hole
\cite{HPSExperiment2022}. Figure~\ref{fig:hps-svt-eras} shows the baseline and
upgraded SVT layouts side by side because that detector evolution is what later drives
the differences in mass resolution and low-mass prompt acceptance used throughout the
analysis note.

\begin{figure}[H]
\centering
\begin{subfigure}[t]{0.90\linewidth}
  \centering
  \graphicorplaceholder{\linewidth}{apparatus_figs/hps_2015_2016_svt_baseline.png}
  \caption{Baseline six-layer SVT used in the 2015 and 2016 engineering runs
  \cite{HPS2016BumpInternal}.}
\end{subfigure}

\vspace{0.9em}

\begin{subfigure}[t]{0.78\linewidth}
  \centering
  \graphicorplaceholder{\linewidth}{apparatus_figs/hps_2019_2021_svt_upgrade.png}
  \caption{Front-layer upgrade used in the 2019 and 2021 physics runs, including the
  added seventh double layer and slim-edge front sensors
  \cite{HPSExperiment2022,HPSSlimEdgeSensors2020}.}
\end{subfigure}
\caption{SVT configurations relevant for the prompt-resonance note. The 2015/2016
baseline tracker and the 2019/2021 upgraded front tracker differ in both layer count
and the construction of the innermost stations, while the ECal remains the same
subsystem apart from the addition of the positron hodoscope before the physics runs.}
\label{fig:hps-svt-eras}
\end{figure}

\subsection{Data-taking campaigns and present analysis scope}

The present note is organized around three HPS running campaigns: the 2015 engineering
run, the 2016 engineering run, and the 2021 physics run. The 2015 and 2016 datasets
are especially natural to compare because they were recorded with the same baseline
detector configuration and were already analyzed in the published prompt-search
workflow. The 2021 dataset belongs to the upgraded detector configuration and is the
first large-statistics physics-run sample to be carried through the GPR workflow.

The 2019 physics run is acknowledged because it shares the upgraded SVT and hodoscope
configuration with 2021, but no 2019 data enter the prompt-resonance result documented
here.

Table~\ref{tab:dataset-summary} summarizes the run periods, analyzed samples, search
intervals, and GP fit supports used in this review note. The last two columns are kept
separate deliberately: a mass hypothesis is tested only inside the search interval,
whereas the GP is trained on the wider fit support after removing the local
$2.25\,\sigma_m$ exclusion.

\begin{table}[H]
\centering
\small
{\setlength{\tabcolsep}{4pt}
\begin{tabularx}{\linewidth}{@{} P{0.20\linewidth} P{0.12\linewidth} P{0.28\linewidth} P{0.16\linewidth} Y @{}}
\toprule
Run period & Beam energy & Sample used here & Search interval & GP fit support \\
\midrule
2015 engineering run &
\SI{1.056}{GeV} &
\SI{1170}{nb^{-1}} &
\SIrange{19}{90}{MeV} &
\SIrange{14}{135}{MeV} \\
2016 engineering run &
\SI{2.3}{GeV} &
\SI{10608}{nb^{-1}} &
\SIrange{39}{180}{MeV} &
\SIrange{30}{210}{MeV} \\
2021 physics run &
\SI{3.74}{GeV} &
10\% distributed development subset of a $\sim\SI{160}{pb^{-1}}$ parent run &
\SIrange{50}{250}{MeV} &
\SIrange{40}{300}{MeV} \\
\bottomrule
\end{tabularx}
}
\caption{Data-taking campaigns carried through the present prompt-resonance note. The
2015 and 2016 inputs are the full engineering-run samples, while the 2021 input is
the reviewed 10\% histogram. The search intervals are unchanged from the recent
combined candidate. Version 4 widens only the GP fit supports, restoring training
sidebands beyond both endpoints of every search. The 2015 upper and 2016 upper fit
edges follow the earlier functional-form validation domains; the 2016 lower edge is
moved to \SI{30}{MeV} because the historical \SI{35}{MeV} edge leaves no rebinned
training-bin center below the \SI{39}{MeV} exclusion. The corresponding published
prompt searches scanned \SIrange{19}{81}{MeV} in 2015 and
\SIrange{39}{179}{MeV} in 2016 \cite{HPS2015DarkPhoton,HPS2016PromptLong}.}
\label{tab:dataset-summary}
\end{table}

The integrated luminosity and analyzed fraction are not inputs to the GP background
interpolation itself. They enter the dataset-specific signal normalization that
converts an extracted prompt yield into an equivalent $\epsilon^2$ limit, together
with the target, acceptance, efficiency, and radiative-fraction inputs for that run
period.

The staged-analysis logic still matters for interpretation. The 2015 and 2016 samples
are now fully unblinded inputs to the shared-coupling result. The 2021 input remains a
distributed 10\% development stream rather than a contiguous end-of-run block, so it
samples run-period conditions while reserving most of that exposure for the final
unblinding. Its ROOT file does not contain enough embedded exposure metadata to turn
the histogram-count ratio to the earlier 1\% file into a luminosity measurement; the
10\% label records the dataset construction rather than an inferred luminosity from
event counts. Version 4 is consequently a full-2015, full-2016, and 2021-10\%
combined analysis, not the final statement for the 2021 run.

Figure~\ref{fig:dataset-mass-shape-overlay} compares the unit-normalized invariant-mass
shapes for the three samples. Figure~\ref{fig:dataset-mass-distributions-log} then shows
the same spectra in absolute events per MeV, with the v4 fit supports and signal-search
intervals distinguished explicitly.

\begin{figure}[H]
\centering
\graphicorplaceholder{0.86\linewidth}{dataset_summary_figs/invariant_mass_distributions_2015_2016_2021_normalized.png}
\caption{Unit-normalized invariant-mass shapes for the full 2015 and 2016 samples
and the 2021 10\% sample used in the v4 workflow. Each histogram is divided by its
own total event count and bin width, so the overlay compares spectral shape rather
than sample size.}
\label{fig:dataset-mass-shape-overlay}
\end{figure}

\begin{figure}[H]
\centering
\graphicorplaceholder{0.98\linewidth}{dataset_summary_figs/invariant_mass_distributions_2015_2016_2021_log.png}
\caption{Absolute invariant-mass distributions for the three v4 inputs. Gray
shading marks the fixed GP fit support and gold shading marks the narrower interval
in which signal hypotheses are tested. The mass-dependent
$2.25\,\sigma_m$ exclusion is removed from the GP support separately at every tested
mass.}
\label{fig:dataset-mass-distributions-log}
\end{figure}
